% 分类目录：dl
\documentclass[12pt, a4paper]{article}
\usepackage[margin=2.5cm]{geometry}
\usepackage{ctex}
\usepackage{amsmath, amssymb}
\usepackage{listings}
\usepackage{xcolor}
\usepackage{tabularx}

\lstset{
    language=Python,
    basicstyle=\ttfamily\small,
    keywordstyle=\color{blue},
    commentstyle=\color{green!50!black},
    stringstyle=\color{red},
    showstringspaces=false,
    breaklines=true,
    numbers=left,
    numberstyle=\tiny\color{gray},
    frame=single
}

\title{知识点：多层感知机 (MLP)}
\author{}
\date{}

\begin{document}
\maketitle

\section{核心定义}
\textbf{中文定义}：多层感知机（Multi-Layer Perceptron, MLP），又称人工神经网络（ANN）或全连接网络，是一种前馈神经网络。它由输入层、一个或多个隐藏层以及输出层组成，层与层之间通过全连接（Fully Connected）的方式连接，即上一层的所有神经元都与下一层的所有神经元相连。

\textbf{英文定义}：A Multi-Layer Perceptron (MLP) is a class of feedforward artificial neural network (ANN). It consists of at least three layers of nodes: an input layer, a hidden layer, and an output layer. Except for the input nodes, each node is a neuron that uses a nonlinear activation function.

\textbf{历史地位}：MLP 是现代深度学习的基石。1980年代反向传播（Backpropagation）算法的引入，使得训练深层 MLP 成为可能，从而克服了早期感知机无法解决非线性（如 XOR 问题）的局限。

\section{核心原理与计算过程}

\subsection{1. 层级结构}
\begin{itemize}
    \item \textbf{输入层}：接收原始特征向量 $x \in \mathbb{R}^d$。
    \item \textbf{隐藏层}：负责提取非线性特征。每个神经元执行：$a = \sigma(Wx + b)$。
    \item \textbf{输出层}：根据任务输出结果（如回归用线性，分类用 Softmax）。
\end{itemize}

\subsection{2. 计算公式}
对于第 $l$ 层，其输出 $h^{(l)}$ 的计算公式为：
$$ z^{(l)} = W^{(l)} h^{(l-1)} + b^{(l)} $$
$$ h^{(l)} = \sigma(z^{(l)}) $$
其中 $W^{(l)}$ 是权重矩阵，$b^{(l)}$ 是偏置项，$\sigma$ 是非线性激活函数。

\subsection{3. 通用近似定理 (Universal Approximation Theorem)}
该定理指出：一个拥有至少一个隐藏层（且包含足够多神经元）的前馈神经网络，只要配以合适的非线性激活函数，就可以以任意精度逼近闭区间上的任何连续函数。

\section{模型结构示意图}
\begin{verbatim}
[ Input x1 ] --\ /-- [ Hidden h1 ] --\ /-- [ Output y1 ]
                X                     X
[ Input x2 ] --/ \-- [ Hidden h2 ] --/ \-- [ Output y2 ]
      |                  |                  |
   (d-dim)           (k-dim)             (classes)
\end{verbatim}

\section{关键代码片段 (PyTorch实现)}
\begin{lstlisting}
import torch.nn as nn

class MLP(nn.Module):
    def __init__(self, input_size, hidden_size, num_classes):
        super(MLP, self).__init__()
        # 定义网络层序列
        self.network = nn.Sequential(
            # 第一层全连接 (Linear / Dense layer)
            nn.Linear(input_size, hidden_size),
            nn.ReLU(),
            # 第二层全连接
            nn.Linear(hidden_size, hidden_size),
            nn.ReLU(),
            # 输出层
            nn.Linear(hidden_size, num_classes)
        )

    def forward(self, x):
        # x shape: (batch_size, input_size)
        return self.network(x)
\end{lstlisting}

\section{深度面试题与解答}
\subsection{基础概念}
\textbf{问：为什么加入隐藏层和激活函数后，模型就能解决非线性问题？}\\
答：如果没有隐藏层的非线性激活，多层全连接网络在数学上等价于单层矩阵相乘，即 $W_2(W_1 x) = (W_2 W_1)x$，本质仍是线性模型。激活函数为每一层引入了扭曲空间的能力，多层叠加后可以构建极复杂的非线性决策边界。

\subsection{进阶原理}
\textbf{问：MLP 的参数量如何计算？}\\
答：设输入维度为 $D_i$，输出维度为 $D_o$。一层全连接层的参数量为：$D_i \times D_o (\text{Weights}) + D_o (\text{Biases})$。

\subsection{工程实战}
\textbf{问：MLP 在处理图像、文本时有什么局限性？}\\
答：1. \textbf{参数冗余}：全连接会忽略数据的空间（图像）或时序（文本）结构，导致参数量随输入尺寸爆炸式增长。2. \textbf{缺乏不变性}：无法利用图像的平移不变性，物体在图中移动一个像素，对于 MLP 来说就是完全不同的输入。

\subsection{坑点分析}
\textbf{问：如果隐藏层神经元非常多，会导致什么问题？}\\
答：容易产生\textbf{过拟合}（Overfitting）。模型会不仅学习到数据的真实规律，还会记住训练集的噪声。解决方法包括使用 Dropout、L2 正则化或减小隐藏层维度。

\section{MLP 与其他模型对比表}

\begin{table}[h]
\centering
\begin{tabularx}{\textwidth}{|l|X|}
\hline
\textbf{模型} & \textbf{核心特征对比} \\
\hline
\textbf{感知机} & 单层，无隐藏层，只能解决线性可分问题。 \\
\hline
\textbf{MLP} & 多层，万能近似器，但参数量大，不适合复杂结构数据。 \\
\hline
\textbf{CNN} & 局部连接、权重共享，擅长处理二维网格结构（图像）。 \\
\hline
\textbf{Transformer} & 基于注意力机制，擅长处理长程依赖和海量数据。 \\
\hline
\end{tabularx}
\end{table}

\section{参考文献与拓展}
\begin{enumerate}
    \item \textbf{经典文献}: Rumelhart, D. E., et al. (1986). \textit{Learning representations by back-propagating errors}. Nature.
    \item \textbf{深度学习圣经}: \textit{Deep Learning} (Ian Goodfellow, Chapter 6: Deep Feedforward Networks).
    \item \textbf{在线课程}: Andrew Ng's Neural Networks and Deep Learning (Coursera).
\end{enumerate}

\end{document}
